% =====================================================================
%  Deliverable 2 — Generative Artificial Intelligence (580694)
%  Universidad de Concepción · 30 de septiembre de 2026
%  Joaquín Sepúlveda · Renato Saavedra · Lucas Saldivia
%
%  UNA PÁGINA, VERTICAL, compilado desde LaTeX. Overleaf → pdfLaTeX.
% =====================================================================
\documentclass[9pt,a4paper]{extarticle}

\usepackage[T1]{fontenc}
\usepackage[utf8]{inputenc}
\usepackage[spanish,es-noquoting]{babel}
% babel-spanish vuelve activa la comilla doble: "c se convierte en ç y rompe
% el JSON de ejemplo. Esto la desactiva.
\shorthandoff{"}
\usepackage{libertine}
\usepackage[scaled=0.80]{beramono}
\usepackage{microtype}
\usepackage[margin=9mm]{geometry}
\usepackage{multicol}
\usepackage{xcolor}
\usepackage{booktabs}
\usepackage{tikz}
\usepackage{array}
\usepackage{colortbl}
\usetikzlibrary{arrows.meta,positioning,fit}

\definecolor{tinta}{HTML}{12191A}
\definecolor{acento}{HTML}{B8501C}
\definecolor{verde}{HTML}{146B5F}
\definecolor{ambar}{HTML}{C08A24}
\definecolor{gris}{HTML}{5A6A6B}
\definecolor{fondo}{HTML}{EDF1F0}
\definecolor{linea}{HTML}{C8D4D1}

\color{tinta}
\pagestyle{empty}
\setlength{\columnsep}{5mm}
\setlength{\parindent}{0pt}
\setlength{\parskip}{0.8mm}
\renewcommand{\arraystretch}{0.97}

\newcommand{\seccion}[1]{\vspace{0.9mm}{\bfseries\footnotesize\color{acento}\MakeUppercase{#1}}\par\nobreak\vspace{0.4mm}\nobreak}
\newcommand{\sub}[1]{{\bfseries #1}\par\nobreak\vspace{-1.2mm}\nobreak}
\newcommand{\q}{\textquotedbl}
\newcommand{\cod}[1]{{\ttfamily\small #1}}
\newcommand{\destaque}[1]{%
  \begingroup\setlength{\fboxsep}{2.0mm}%
  \colorbox{fondo}{\parbox{\dimexpr\columnwidth-4.0mm\relax}{#1}}%
  \endgroup\par}

\begin{document}
\raggedright

% ---------------------------------------------------------------- cabecera
{\Large\bfseries De 16,7\,\% a 58,3\,\% en la validación de inscripción de ramos con un modelo de 3,8B}\\[0.8mm]
{\footnotesize\color{gris}%
\textbf{Deliverable 2} · Generative Artificial Intelligence (580694) · Universidad de Concepción · 30 de septiembre de 2026\\
\textbf{Equipo:} Joaquín Sepúlveda · Renato Saavedra · Lucas Saldivia \quad\textbullet\quad
\textbf{Modelo:} \cod{microsoft/Phi-3.5-mini-instruct}, 3,8B, cuantizado a 4 bits, T4 de Colab\\
\textbf{Repositorio:} \cod{github.com/Luquitas02/genai-inscripcion-ramos} \quad\textbullet\quad
\textbf{Video:} \cod{<PEGAR LINK>}}

\vspace{1.5mm}
{\color{linea}\rule{\textwidth}{0.4pt}}
\vspace{1mm}

\begin{multicols}{2}

% ================================================================ COLUMNA 1
\seccion{La tarea y la falla que se ataca}

Dado el historial de un estudiante y una pregunta escrita como la escribiría una persona,
decidir si puede inscribir un ramo y citar la regla que lo sostiene. Salida de dos campos,
corrección por coincidencia exacta, 60 casos balanceados. Es la misma tarea del Deliverable~1,
con el mismo conjunto y el mismo criterio.

El Deliverable~1 midió que el baseline no razona. Con prompting directo, Phi acierta 16,7\,\%
de los casos en ambos campos, bajo el piso de 30,0\,\% que consigue quien responde siempre lo
mismo. Se diagnosticaron cuatro mecanismos: no representa el estado \q cursándola ahora\q\ (E1),
depende de que la pregunta le apunte al dato (E2), inventa identificadores (E3) y se contradice
entre decisión y regla (E4).

\seccion{La intervención}

Descomposición en tres pasos con contexto enfocado y salida restringida. El baseline recibe
3.862 tokens y debe hacer todo en una pasada. Acá cada paso recibe solo lo suyo.

\vspace{1mm}
\begin{center}
\begin{tikzpicture}[
  font=\scriptsize,
  caja/.style={draw=linea,fill=white,rounded corners=1pt,inner sep=1.4mm,text width=30mm,align=center},
  modelo/.style={caja,draw=acento,line width=0.6pt},
  codigo/.style={caja,draw=verde,line width=0.6pt},
  fl/.style={-{Stealth[length=1.6mm]},draw=gris}]

\node[caja,fill=fondo] (q) {\q me falta cálculo 2 pero la estoy\\ tomando, ¿alcanzo a meter\\ ecuaciones el próximo?\q};
\node[modelo,below=2.6mm of q] (p1) {\textbf{Paso 1 · Extracción}\\[0.3mm] {\color{acento}el modelo}\\ pregunta + 61 códigos};
\node[codigo,below=2.6mm of p1] (p2) {\textbf{Paso 2 · Ficha}\\[0.3mm] {\color{verde}el código}\\ 8 campos, 1 ramo};
\node[codigo,below=2.6mm of p2] (p3) {\textbf{Paso 3 · Decisión}\\[0.3mm] {\color{verde}el código}\\ lista cerrada de reglas};
% La comilla doble literal la come babel-spanish aun con \shorthandoff. Va \q, que es
% \textquotedbl y es inmune, igual que en el resto del documento.
\node[caja,fill=fondo,below=2.6mm of p3] (out) {\cod{\{\q regla\q: \q R-DEPENDE-APROBACION\q\}}\\ decisión derivada: \textbf{condicional}};

\draw[fl] (q) -- (p1);
\draw[fl] (p1) -- node[right=0.5mm,font=\tiny,text=gris]{ramo, período} (p2);
\draw[fl] (p2) -- node[right=0.5mm,font=\tiny,text=gris]{descarta 25 asignaturas} (p3);
\draw[fl] (p3) -- (out);
\end{tikzpicture}
\end{center}
\vspace{0.5mm}

{\footnotesize
El \textbf{paso 1} ataca E2. Ve la pregunta y los 61 pares de código y nombre, sin historial ni
reglas. El \textbf{paso 2} ataca E1. Arma una ficha de ocho campos sobre un solo ramo y descarta
las 25 asignaturas del historial que no lo tocan. El \textbf{paso 3} ataca E3 y E4. La regla
debe salir de una lista cerrada que incluye los códigos de los prerrequisitos del caso, y la
decisión se deriva de esa regla.
}

\seccion{Dos mecanismos eliminados por construcción}

Se verificó sobre los 60 casos que la decisión queda determinada por la regla, sin una sola
ambigüedad. Pedir los dos campos le daba al modelo un grado de libertad que el dominio no tiene.
Al derivar la decisión y restringir la regla, E3 y E4 desaparecen.

\vspace{0.5mm}
{\footnotesize
\begin{tabular}{@{}lr@{}}
\toprule
sobre 240 respuestas donde el modelo decide & \\
\midrule
identificadores inventados (E3) & \textbf{0} \\
contradicciones decisión/regla (E4) & \textbf{0} \\
salidas sin JSON legible & \textbf{0} \\
\bottomrule
\end{tabular}}

\vspace{1mm}
En el baseline zero-shot, Phi inventaba identificadores en 18 de 60 casos.

\seccion{Acierto por tipo de caso}

Los 60 casos se generaron desde la regla que debía decidir, así que cada tipo se puede
contar por separado. El sistema mejora en siete de las ocho familias.

\vspace{0.5mm}
{\footnotesize
\begin{tabular}{@{}lrr@{}}
\toprule
regla que decide el caso & baseline & sistema \\
\midrule
sin impedimento              & 2/18 & \textbf{13/18} \\
prerrequisito en curso, próximo período & 6/12 & 3/12 \\
excepción por carga baja     & 0/9  & \textbf{8/9} \\
prerrequisito faltante       & 2/6  & \textbf{3/6} \\
prerrequisito en curso, período actual & 0/6 & \textbf{1/6} \\
umbral de créditos aprobados & 0/3  & \textbf{3/3} \\
tope máximo de créditos      & 0/3  & \textbf{2/3} \\
requisito de año completo    & 0/3  & \textbf{2/3} \\
\bottomrule
\end{tabular}}

% ================================================================ COLUMNA 2
\columnbreak

\seccion{Resultados sobre los 60 casos}

Las variantes mueven una pieza por vez del modelo al código, así que la diferencia entre dos
filas mide lo que esa pieza aporta.

\vspace{0.5mm}
{\footnotesize
\begin{tabular}{@{}lrrrr@{}}
\toprule
 & decisión & regla & \textbf{ambas} & tok/caso \\
\midrule
baseline few-shot       & 35,0 & 23,3 & 16,7 & 3.862 \\
constante trivial       & 35,0 & 30,0 & 30,0 & --- \\
\midrule
ficha y reglas al modelo   & 45,0 & 25,0 & 25,0 & 6.072 \\
ficha al código            & 45,0 & 28,3 & 28,3 & 2.743 \\
\rowcolor{fondo}
\textbf{ficha y reglas al código} & \textbf{71,7} & \textbf{58,3} & \textbf{58,3} & \textbf{1.153} \\
\midrule
{\itshape con el paso 1 resuelto} & 100,0 & 100,0 & 100,0 & 1.153 \\
\bottomrule
\end{tabular}}

\vspace{1mm}
\destaque{\footnotesize El sistema triplica el baseline y es además la condición más barata.
Usa \textbf{1.153 tokens por caso contra 3.862} del baseline, y 1,6 segundos contra 4,2. Con un
tercio del cómputo, la mejora no se explica por haber gastado más.}

\seccion{Estrategias alternativas evaluadas}

El paso 3 se midió con la ficha verdadera de cada caso, bajo cuatro versiones del prompt, contra
un techo demostrado de 60/60.

\vspace{0.5mm}
{\footnotesize
\begin{tabular}{@{}lrl@{}}
\toprule
versión del prompt del paso 3 & ambas & regla dominante \\
\midrule
completo                     & 26,7 & \cod{R-EXCEPCION} 38 \\
sin el bloque de reglas      & \textbf{31,7} & \cod{R-DEPENDE} 49 \\
sin ejemplos                 & 15,0 & \cod{R-EXCEPCION} 60 \\
mínimo                       & 30,0 & \cod{R-DEPENDE} 47 \\
\bottomrule
\end{tabular}}

\vspace{1mm}
{\footnotesize
El bloque de reglas globales dictaba la respuesta por su posición. Era la última prosa del
prompt y la única que nombraba un valor de decisión. Sacarlo vale 5 puntos. Aun así el modelo
sigue colapsando en una categoría. Con la ficha perfecta delante aplica las reglas al 31,7\,\%
y el código lo hace al 100\,\%. Por esa diferencia las reglas quedaron en el código.

En el paso 1 los ejemplos resueltos se midieron y se descartaron. Bajan el acierto de ramo de
51,7\,\% a 45,0\,\%.}

\seccion{Dónde se rompe, y por qué}

\vspace{0.5mm}
{\footnotesize
\begin{tabular}{@{}lr@{}}
\toprule
casos donde el paso 1 acierta ramo y período & \textbf{28 de 28} \\
casos donde el paso 1 falla & 7 de 32 \\
\bottomrule
\end{tabular}}

\vspace{1mm}
{\footnotesize
\textbf{Todo el error restante es extracción.} El paso 1 identifica el ramo en 31 de 60 casos, y
se degrada con la forma de nombrarlo: 13/20 con el nombre completo, 10/20 abreviado, 8/20
coloquial. Las 29 fallas devolvieron un código válido pero equivocado, ninguna devolvió basura.

\textbf{Caso de falla, elegido por regla escrita antes de medir} (el primero de nivel 3 cuyo
baseline falla). La pregunta es \q programación, la reprobé, ¿puedo tomar Optimización 1 ahora
ya?\q. El ramo consultado es 580315. El paso 1 devolvió 503203, que la frase menciona de
contexto y que además es el prerrequisito que bloquea. Los pasos 2 y 3 evaluaron 503203 y
respondieron que sí, porque un ramo reprobado se puede volver a inscribir. La respuesta es
correcta para ese ramo y ese ramo no era el consultado. Es el mecanismo E2 del Deliverable~1.

\textbf{El paso 2 tampoco puede leer el expediente} aun con el ramo correcto: estado del ramo
0/60, créditos aprobados 0/60, prerrequisitos 7/60. Por eso la ficha la arma el código.}

\seccion{Límites declarados}

{\footnotesize
\textbf{Una categoría empeora.} En los 12 casos con prerrequisito en curso preguntando por el
próximo período, el sistema baja de 6/12 a 3/12. Los 12 tienen la misma respuesta correcta, y
el baseline la dio en 6 de ellos respondiendo \q condicional\q. Las 9 fallas del sistema son
las 9 veces que el paso 1 tomó el ramo equivocado.

\textbf{Filtración del conjunto.} Los 9 casos de \cod{R-EXCEPCION-PRERREQ} tienen todos cero
créditos inscritos y los otros 51 tienen entre 12 y 24, con el umbral en 8. El conjunto no
prueba la frontera de esa regla.

\textbf{Selección sobre el mismo conjunto.} Las versiones del prompt se eligieron por ablación
sobre los mismos 60 casos. Con n=60, diferencias de dos a cuatro casos no deben leerse como
efectos firmes.

\textbf{Desviación declarada.} El plan del Deliverable~1 reservaba la consulta al grafo para el
Deliverable~3. Se adelantó porque la ablación mostró que el modelo no aplica las reglas. El
Deliverable~3 pasa a atacar la extracción, que es el cuello de botella medido.

\textbf{Un dominio, 60 casos.} Sin cambios respecto al Deliverable~1.}

\end{multicols}
\end{document}
